% Pass7509--7516: identify the common S4 of the Monster-local 9+3360 action.
\subsection*{Passes 7509--7516: the common $S_4$ is the affine-direction quotient}

Pass7501--7508 identified the Monster maximal subgroup
\[
(3^2{:}2\times O_8^+(3)).S_4
\cong
AGL(2,3)\times_{S_4}(O_8^+(3):S_4)
\]
from the hash-pinned ATLAS $3369$-point action.  The common $S_4$ can now be
identified geometrically rather than merely by its order.

The restricted nine-point image has order $432$ and contains a normal regular
translation subgroup
\[
T\cong C_3^2.
\]
Its four order-three subgroups are the four affine directions.  Their point
orbits produce exactly twelve three-point lines in four parallel classes of
three lines each.  Every pair of points lies on one line and every point lies on
four lines, so the recovered incidence geometry is
\[
\boxed{AG(2,3)=(9_4,12_3).}
\]
Conjugation by the full affine group permutes the four direction subgroups.  The
induced image has order $24$, with element-order census
\[
1^1\,2^9\,3^8\,4^6,
\]
and its kernel has order $18$.  Therefore
\[
\boxed{AGL(2,3)/(3^2{:}2)\cong PGL(2,3)\cong S_4.}
\]
Thus the same four-letter coordinate has two exact meanings in the Monster-local
fibre product: on the nine-sheet it labels the four affine/Hesse parallel
classes, while on the $3360$-sheet it is the outer $D_4(3)$ coordinate acting on
the triality geometry.  The sporadic subgroup couples these two direction
systems through this common $S_4$.

\paragraph{Evidence boundary.}
The project already contained the abstract $AG(2,3)$/Hesse counts.  The new
statement is the exact welding of their direction quotient to the outer
$S_4$ of the Monster-local $3360$ triality carrier.  No physical interpretation
of the four directions follows.
